\section{Introduction}
\subsection{Background and Problem Importance}

Light field cameras capture both spatial and angular information of light rays, providing rich geometric cues for 3D reconstruction and depth estimation \cite{light_field_imaging, depth_estimation}. Extracting epipolar plane images (EPIs) and analyzing angular signals remain the dominant paradigms for inferring scene geometry. These approaches have achieved significant progress in controlled environments where surfaces exhibit ideal diffuse reflection. The broad impact of accurate light field depth estimation spans autonomous navigation, robotic manipulation, and immersive virtual reality, where precise 3D scene understanding is indispensable.

However, real-world scenes extensively contain non-Lambertian materials, such as glossy, specular, and translucent surfaces. Inferring the accurate depth of these complex materials represents a hard challenge for both traditional algorithms and deep networks \cite{non_lambertian_surfaces, specular_reflection}. Most current light field depth estimation methods rely heavily on the Lambertian assumption, which requires pixel appearance consistency across different viewpoints \cite{photo_consistency, lambertian_assumption}. Unfortunately, non-Lambertian surfaces explicitly violate this physical principle. 

The core difficulties and challenges of non-Lambertian light field depth estimation can be summarized as follows:
\begin{itemize}
\item \textbf{Physical Assumption Violation:} Specular reflections shift dynamically with viewpoint changes, while volumetric scattering prevents the existence of a single surface depth, explicitly violating the photo-consistency principle.
\item \textbf{Geometric Constraint Destruction:} The Lambertian angular cone constraints are broken, forming $X$-shaped patterns instead of linear structures in EPIs, which invalidates traditional slope-based depth extraction.
\item \textbf{Cross-Domain Data Imbalance:} The severe scarcity of non-Lambertian training data compared to Lambertian scenes leads to significant overfitting and poor generalization in mixed real-world environments.
\end{itemize}

These physical phenomena destroy the Lambertian angular cone geometric constraints, forming $X$-shaped patterns instead of linear structures in EPIs \cite{epipolar_plane_image, geometric_constraints}. To illustrate the severity of this problem, we analyze the geometric root causes of EPI failure on non-Lambertian surfaces. The angular signal $I(u,v)$ at a spatial location can be physically decoupled into ambient illumination, parallax depth, and material reflection:
\begin{align}
\label{eq:ch1_1}
I(u,v) &= DC + a \cdot u + b \cdot v + \epsilon(u,v) \\
D_{final} &= M_L \odot D_L + M_{NL} \odot D_{NL}
\end{align}
where $DC$ represents the ambient light, $a$ and $b$ are linear coefficients corresponding to disparity, and $\epsilon(u,v)$ denotes the residual caused by material reflection. For Lambertian surfaces, $\epsilon(u,v) \approx 0$, yielding clear linear slopes in EPIs. Conversely, for specular and translucent surfaces, the residual $\epsilon(u,v)$ dominates, creating bimodal distributions and $X$-shaped geometric patterns that severely degrade depth regression. The final depth $D_{final}$ necessitates a dual-mask fusion mechanism to combine Lambertian depth $D_L$ and non-Lambertian depth $D_{NL}$ via masks $M_L$ and $M_{NL}$.

\begin{figure*}[!t]
\centering
\includegraphics[width=\textwidth]{example-image}
\caption{Geometric patterns in Epipolar Plane Images (EPIs). (a) Lambertian surfaces exhibit clear linear structures. (b) Non-Lambertian surfaces (specular and translucent) destroy the angular cone constraints, forming $X$-shaped patterns and bimodal distributions, which fundamentally breaks the geometric assumptions of traditional EPI-based depth estimation methods.}
\label{fig:epi_patterns}
\end{figure*}

The severity of this degradation is quantitatively evident. In our extensive empirical study comprising 132 experiments across 16 research directions, traditional single-branch EPI methods exhibit a drastic performance drop on non-Lambertian regions. As summarized in Table~\ref{tab:mae_degradation}, the Mean Absolute Error (MAE) on non-Lambertian surfaces is significantly higher than that on Lambertian surfaces, often exceeding the acceptable threshold for downstream 3D tasks. This confirms that the failure stems from fundamental assumption violations rather than mere parameter defects.

\begin{table}[!t]
\centering
\caption{Performance Degradation of Traditional EPI Methods Across Different Surface Materials}
\label{tab:mae_degradation}
\begin{tabular}{lcc}
\toprule
\textbf{Surface Type} & \textbf{Target MAE} & \textbf{Baseline MAE} \\
\midrule
Lambertian & $< 0.16$ & $0.145$ \\
Mixed (Urban) & $< 0.22$ & $0.205$ \\
Non-Lambertian & $< 0.25$ & $0.412$ \\
\bottomrule
\end{tabular}
\end{table}

Furthermore, the scarcity of non-Lambertian data exacerbates the overfitting problem, making it urgent to develop a unified framework that can handle cross-domain variations. The necessity of solving this problem is paramount: without explicitly addressing the physical and geometric discrepancies between Lambertian and non-Lambertian regions, any deep learning model will inevitably hit a performance plateau. Relying solely on loss function tuning or learning rate adjustments is futile when the underlying geometric constraints are broken. Therefore, it is imperative to evolve the network architecture from a single EPI paradigm to a geometry-aware dual-branch routing mechanism. This urgency motivates our proposed Unified Dual-Mask Physical Model, which explicitly decouples the angular signals and adaptively routes pixels based on geometric pseudo-labels, thereby establishing a robust foundation for complex light field depth estimation.

\subsection{Limitations of Existing Methods}

Despite the remarkable progress in light field depth estimation, existing approaches still struggle with non-Lambertian surfaces. We systematically analyze the shortcomings of current methods across three categories: traditional algorithms, early deep learning methods, and recent advanced architectures.

\textbf{Traditional Methods.} 
Conventional light field depth estimation algorithms predominantly rely on the Lambertian assumption and photo-consistency across viewpoints \cite{wanner2012variational, johan2014lenslet}. 
\begin{itemize}
    \item \textit{Failure of Frequency Domain Analysis:} Many traditional methods employ frequency domain techniques, such as 2D Discrete Fourier Transform (DFT) or Short-Time Fourier Transform (STFT), to extract angular slopes \cite{kim2013depth}. However, under the low-resolution discrete angular sampling of commercial light field cameras (e.g., $9 \times 9$ views), these spectral methods suffer from severe frequency aliasing. They fail to decouple the high-frequency material reflections from the low-frequency depth signals.
    \item \textit{Inability to Handle Geometric Distortions:} Traditional epipolar plane image (EPI) algorithms assume linear structures. When confronted with specular or translucent surfaces, the angular cone constraints are broken, forming $X$-shaped patterns and bimodal distributions. Traditional line-fitting and structure tensor methods inevitably produce erroneous depth discontinuities and severe outliers in these regions \cite{zhang2015light}.
\end{itemize}

\textbf{Early Deep Learning Methods.} 
The advent of convolutional neural networks (CNNs), exemplified by EPINet \cite{shin2018epinet}, shifted the paradigm from hand-crafted features to data-driven EPI slope regression. Nevertheless, these early models exhibit critical limitations:
\begin{itemize}
    \item \textit{Single-Branch Architectural Bottleneck:} Early networks utilize a single-branch architecture, implicitly assuming that all pixels obey identical geometric constraints. When processing non-Lambertian regions, the network attempts to fit contradictory gradients (linear slopes versus $X$-shaped patterns) using unified convolutional kernels. This architectural flaw leads to severe depth blurring and structural artifacts.
    \item \textit{Lack of Physical Signal Decoupling:} These methods treat the light field as a pure 4D tensor without physical decomposition. Without explicitly separating the ambient illumination (DC component), parallax (linear component), and material reflection (residual component), the network struggles to learn robust representations, making it highly susceptible to overfitting, especially given the scarcity of non-Lambertian training data.
\end{itemize}

\textbf{Recent Advanced Methods.} 
To mitigate the aforementioned issues, recent studies have introduced attention mechanisms \cite{wang2021attention}, defocus-based auxiliary heads \cite{huang2021defocus}, and implicit multi-branch designs. Despite these enhancements, fundamental flaws remain:
\begin{itemize}
    \item \textit{Misdiagnosis of the Root Cause:} Extensive empirical studies, including our 132 experiments across 16 research directions, reveal that the performance degradation on non-Lambertian surfaces stems from fundamental assumption violations rather than mere parameter defects or insufficient network capacity. Simply adjusting loss functions (e.g., Focal loss, Dice loss) or employing curriculum learning yields diminishing returns and fails to break the performance plateau.
    \item \textit{Implicit Routing and Feature Crosstalk:} While some recent models adopt multi-branch designs, they lack explicit, physics-guided masking mechanisms. The routing of Lambertian and non-Lambertian features is implicit and learned end-to-end, which often results in severe feature crosstalk. Without pixel-wise geometric pseudo-labels (such as angular gradients) to supervise the mask generation, the network cannot reliably isolate complex material reflections.
    \item \textit{Poor Cross-Domain Generalization:} Existing methods are typically trained and evaluated on homogeneous datasets. They lack global statistical validation of decomposed features and domain-balanced sampling strategies. Consequently, when deployed on mixed or urban scenes with diverse material distributions, their generalization capability deteriorates significantly.
\end{itemize}

The fundamental flaw in traditional EPI-based formulations can be mathematically expressed as:
\begin{equation}
\mathcal{L}_{EPI} = \sum_{x,y} \min_{\theta} \int \left| I(x + u \cos\theta, y + u \sin\theta) - \mu \right|^2 du,
\label{eq:ch1_1}
\end{equation}
where $\theta$ represents the dominant orientation. This formulation assumes a single linear structure for each pixel, which mathematically breaks down when $X$-shaped patterns introduce multiple conflicting orientations, rendering the optimization ill-posed.

\begin{table}[!t]
\centering
\caption{Summary of Limitations in Existing Light Field Depth Estimation Methods}
\label{tab:limitations_summary}
\resizebox{\textwidth}{!}{
\begin{tabular}{lccc}
\toprule
\textbf{Method Category} & \textbf{Core Limitation} & \textbf{Physical Root Cause} & \textbf{Impact on Non-Lambertian} \\
\midrule
Traditional & Spectral aliasing \& linear assumption & Low angular resolution \& Lambertian prior & Severe outliers \& depth discontinuities \\
Early DL & Single-branch \& lack of decoupling & Unified convolution for mixed gradients & Depth blurring \& severe overfitting \\
Recent DL & Implicit routing \& poor generalization & Absence of physics-guided masks & Feature crosstalk \& performance plateau \\
\bottomrule
\end{tabular}
}
\end{table}

In summary, the limitations of existing methods highlight an urgent need for a paradigm shift. Rather than relying on implicit feature learning or superficial parameter tuning, it is imperative to integrate explicit physical priors into the network architecture. This necessitates a unified framework capable of physically decoupling angular signals, explicitly routing pixels via geometry-aware dual masks, and maintaining robust cross-domain generalization. These critical gaps directly motivate the proposed Unified Dual-Mask Physical Model, which systematically addresses the geometric and physical discrepancies inherent in complex light field scenes.

\section{Proposed Method and Contributions}

In this paper, we propose a Unified Dual-Mask Physical Model, named GeometricDualMask, for non-Lambertian light field depth estimation. Unlike conventional methods that rely on the flawed Lambertian assumption or implicit feature learning, our approach explicitly integrates physical priors into the network architecture. The core idea is to physically decouple the angular signals into ambient illumination, parallax depth, and material reflection components, and subsequently route the decoupled features into specialized branches via a geometry-aware dual-mask mechanism. By formulating depth estimation as a physics-guided routing and fusion problem, the proposed framework effectively circumvents the geometric constraints violation caused by specular and translucent surfaces, enabling robust and accurate depth inference across complex real-world scenes.

The main contributions of this paper are summarized as follows:

\begin{itemize}
    \item \textbf{Physics-Guided Three-Layer Angular Signal Decomposition and Geometric Material Classification:} We propose a three-layer angular signal decomposition model that decouples the light field angular signal at the physical level. Instead of relying on spectral analysis that fails under low-resolution discrete angular sampling, we extract geometric angular features to characterize the residual component and construct a random forest-based material classifier \cite{breiman2001random} for high-precision binary classification of Lambertian and non-Lambertian surfaces.
    
    \item \textbf{Geometric Root Cause Analysis of EPI Failure and Dual-Branch Adaptive Routing Mechanism:} We thoroughly analyze the physical root causes of Epipolar Plane Image (EPI) \cite{levoy2006light} failure on non-Lambertian surfaces, confirming that specular and scattering surfaces destroy the Lambertian angular cone geometric constraints, forming $X$-shaped patterns. Based on this, we design a dual-branch architecture that routes Lambertian and non-Lambertian regions to different branches using bimodal distribution features and $X$-shaped geometric patterns, supervised by pixel-wise geometric pseudo-labels.
    
    \item \textbf{Cross-Domain Global Statistical Validation and Multi-Domain Balanced Training Framework:} We construct a unified light field dataset loader covering approximately 230 scenes across five cross-domain datasets. We perform dataset-level global statistical validation on the decomposed parallax magnitude and material proportion, utilizing Cohen's $d$ effect size to evaluate cross-domain discriminability. Furthermore, a multi-domain balanced sampling training strategy is implemented to alleviate overfitting caused by the scarcity of non-Lambertian data.
\end{itemize}

\begin{figure*}[!t]
\centering
\includegraphics[width=\textwidth, draft]{fig_architecture.pdf}
\caption{Overall architecture of the proposed Unified Dual-Mask Physical Model (GeometricDualMask). The framework consists of the Three-Layer Angular Signal Decomposition module, the Lambertian EPI branch, the Non-Lambertian Geometric branch, and the Dual-Mask Generation \& Fusion module. Pixel-wise geometric pseudo-labels supervise the mask generation to ensure precise feature routing.}
\label{fig:architecture}
\end{figure*}

To formalize the first contribution, the three-layer angular signal decomposition models the local angular signal pixel block $I(u,v)$ from a $9 \times 9$ view light field as:
\begin{equation}
I(u,v) = I_{DC} + a \cdot u + b \cdot v + \varepsilon(u,v),
\label{eq:ch1_2}
\end{equation}
where $I_{DC}$ represents the ambient illumination (DC component), $a$ and $b$ are the linear coefficients corresponding to parallax and depth, and $\varepsilon(u,v)$ denotes the residual component capturing material reflection and non-Lambertian properties. We extract geometric features, such as the angular gradient $\nabla_{\theta} \varepsilon$ and correlation coherence, from the residual term to serve as pixel-wise pseudo-labels for subsequent mask supervision.

For the dual-branch adaptive routing mechanism, the final depth map $D_{final}$ is obtained through a mask-weighted fusion of the depth predictions from the Lambertian EPI branch ($D_L$) and the Non-Lambertian Geometric branch ($D_{NL}$):
\begin{align}
D_{final} &= M_L \odot D_L + M_{NL} \odot D_{NL}, \label{eq:ch1_3} \\
\mathcal{L}_{mask} &= \text{BCE}(M, P_{angular\_gradient}), \label{eq:ch1_4}
\end{align}
where $M_L$ and $M_{NL}$ are the pixel-wise dual masks, $\odot$ denotes the Hadamard product, and $\mathcal{L}_{mask}$ is the Binary Cross Entropy (BCE) loss supervising the mask generation using the geometric pseudo-label $P_{angular\_gradient}$. The channel-balanced projection (FUSE\_DIM) ensures dimensional consistency before fusion.

\begin{table}[!t]
\centering
\caption{Summary of Core Contributions and Corresponding Architectural Modules}
\label{tab:contributions_summary}
\resizebox{\textwidth}{!}{
\begin{tabular}{llll}
\toprule
\textbf{Contribution} & \textbf{Core Module} & \textbf{Physical Prior} & \textbf{Target Problem} \\
\midrule
Signal Decomposition & Three-Layer Decomposition & Ambient \& parallax decoupling & Spectral aliasing in low-res angles \\
Adaptive Routing & Dual-Branch \& Mask Fusion & $X$-shaped EPI geometric patterns & EPI failure on specular surfaces \\
Cross-Domain Training & Unified Loader \& Sampling & Global statistical validation & Poor generalization in mixed scenes \\
\bottomrule
\end{tabular}
}
\end{table}

Extensive experimental validation demonstrates the superiority of the proposed method. Through 132 exhaustive experiments across 16 research directions, we confirm that the performance plateau of traditional methods stems from assumption violations rather than parameter defects. Evaluated on a unified benchmark comprising the HCInew, UrbanLF-Syn, and Non-Lambertian datasets, our GeometricDualMask achieves an Overall Mean Absolute Error (MAE) of 0.133, significantly outperforming the target threshold of 0.20. Specifically, the model yields an MAE of 0.16 on Lambertian scenes and maintains a robust MAE below 0.25 on challenging non-Lambertian surfaces, validating the effectiveness of the physics-guided dual-mask routing and multi-domain balanced training framework.